\documentclass[11pt]{article}

\usepackage[margin=1in]{geometry}
\usepackage{amsmath}
\usepackage{amssymb}
\usepackage{booktabs}
\usepackage{array}
\usepackage[hidelinks]{hyperref}

\newcommand{\softmax}{\operatorname{softmax}}

\title{\textbf{Exercise 05: KV Cache for Autoregressive Decoding}}
\author{}
\date{}

\begin{document}
\maketitle

\begin{center}
\emph{A numerical comparison of full-prefix causal recomputation and cached single-token decoding.}
\end{center}

\section*{Setup}

The prefill is
\[
X_{\mathrm{past}}=
\begin{bmatrix}
1&0\\
0&1\\
1&1
\end{bmatrix},
\]
and the new decode token is
\[
x_t=x_{\mathrm{today}}=\begin{bmatrix}0&1\end{bmatrix}.
\]
We use
\[
W_Q=\begin{bmatrix}1&0\\0&1\end{bmatrix},\qquad
W_K=\begin{bmatrix}1&1\\0&1\end{bmatrix},\qquad
W_V=\begin{bmatrix}1&0\\1&1\end{bmatrix},
\]
with $d_k=d_v=2$.

\section*{Part A --- Prefill}

The projected prefill tensors are
\[
Q_{\mathrm{past}}=X_{\mathrm{past}}W_Q=
\begin{bmatrix}
1&0\\0&1\\1&1
\end{bmatrix},
\]
\[
K_{\mathrm{past}}=X_{\mathrm{past}}W_K=
\begin{bmatrix}
1&1\\0&1\\1&2
\end{bmatrix},
\qquad
V_{\mathrm{past}}=X_{\mathrm{past}}W_V=
\begin{bmatrix}
1&0\\1&1\\2&1
\end{bmatrix}.
\]

For future decode steps, the reusable state is
\[
K_{\mathrm{cache}}=K_{\mathrm{past}},\qquad
V_{\mathrm{cache}}=V_{\mathrm{past}}.
\]
Past query rows are not part of the cache because the next token needs a new query, not the old queries.

\section*{Part B --- No-cache full-prefix causal decode}

Append the new token:
\[
X_{\mathrm{full}}=
\begin{bmatrix}
1&0\\0&1\\1&1\\0&1
\end{bmatrix}.
\]
Recomputing the projections over the full prefix gives
\[
Q_{\mathrm{full}}=
\begin{bmatrix}
1&0\\0&1\\1&1\\0&1
\end{bmatrix},
\quad
K_{\mathrm{full}}=
\begin{bmatrix}
1&1\\0&1\\1&2\\0&1
\end{bmatrix},
\quad
V_{\mathrm{full}}=
\begin{bmatrix}
1&0\\1&1\\2&1\\1&1
\end{bmatrix}.
\]

The first three K/V rows are unchanged from prefill, so recomputing them is redundant.

The full scaled score matrix is
\[
S_{\mathrm{full}}
=\frac{Q_{\mathrm{full}}K_{\mathrm{full}}^\top}{\sqrt{2}}
\approx
\begin{bmatrix}
0.7071&0&0.7071&0\\
0.7071&0.7071&1.4142&0.7071\\
1.4142&0.7071&2.1213&0.7071\\
0.7071&0.7071&1.4142&0.7071
\end{bmatrix}.
\]

The causal mask is
\[
M=
\begin{bmatrix}
0&-\infty&-\infty&-\infty\\
0&0&-\infty&-\infty\\
0&0&0&-\infty\\
0&0&0&0
\end{bmatrix}.
\]
Applying row-wise softmax to $S_{\mathrm{full}}+M$ gives
\[
A_{\mathrm{full}}
\approx
\begin{bmatrix}
1.0000&0&0&0\\
0.5000&0.5000&0&0\\
0.2840&0.1400&0.5760&0\\
0.1989&0.1989&0.4034&0.1989
\end{bmatrix}.
\]

Therefore
\[
O_{\mathrm{full}}=A_{\mathrm{full}}V_{\mathrm{full}}
\approx
\begin{bmatrix}
1.0000&0.0000\\
1.0000&0.5000\\
1.5760&0.7160\\
1.4034&0.8011
\end{bmatrix}.
\]
The newest-token output from the no-cache reference is the last row:
\[
\boxed{o_t^{(\mathrm{no-cache})}\approx\begin{bmatrix}1.4034&0.8011\end{bmatrix}}.
\]

\section*{Part C --- Cached decode}

Only the new token is projected:
\[
q_t=x_tW_Q=\begin{bmatrix}0&1\end{bmatrix},
\]
\[
k_t=x_tW_K=\begin{bmatrix}0&1\end{bmatrix},
\qquad
v_t=x_tW_V=\begin{bmatrix}1&1\end{bmatrix}.
\]
Append the new key and value:
\[
K_{\mathrm{cache}}=
\begin{bmatrix}
1&1\\0&1\\1&2\\0&1
\end{bmatrix},
\qquad
V_{\mathrm{cache}}=
\begin{bmatrix}
1&0\\1&1\\2&1\\1&1
\end{bmatrix}.
\]
These are exactly $K_{\mathrm{full}}$ and $V_{\mathrm{full}}$, but the old rows were reused rather than recomputed.

The cached score vector is
\[
s_t^{(\mathrm{cache})}
=\frac{q_tK_{\mathrm{cache}}^\top}{\sqrt{2}}
\approx
\begin{bmatrix}
0.7071&0.7071&1.4142&0.7071
\end{bmatrix}.
\]
Its softmax is
\[
a_t^{(\mathrm{cache})}
\approx
\begin{bmatrix}
0.1989&0.1989&0.4034&0.1989
\end{bmatrix}.
\]
Thus
\[
o_t^{(\mathrm{cache})}
=a_t^{(\mathrm{cache})}V_{\mathrm{cache}}
\approx
\boxed{\begin{bmatrix}1.4034&0.8011\end{bmatrix}}.
\]

\section*{Part D --- Why the outputs are identical}

The newest row of the full causal attention matrix is
\[
A_{\mathrm{full}}[-1]
=\begin{bmatrix}0.1989&0.1989&0.4034&0.1989\end{bmatrix},
\]
which is exactly $a_t^{(\mathrm{cache})}$.

This happens because appending a token does not change the already-computed keys and values of earlier tokens. In a causal model, the newest token needs only:
\begin{itemize}
    \item its current query $q_t$,
    \item all available keys,
    \item all corresponding values.
\end{itemize}
The old query rows were useful when the old positions produced their own outputs, but they are not needed to produce the newest token's output.

Therefore
\[
\boxed{o_t^{(\mathrm{no-cache})}=o_t^{(\mathrm{cache})}}.
\]

\section*{Part E --- Repeated work}

For three decode calls after the common three-token prefill, the prefix lengths are $4$, $5$, and $6$.

\subsection*{Projection rows}
Without a cache, projecting the full prefix each time processes
\[
4+5+6=15
\]
token rows. With a cache, only the three new tokens are projected:
\[
1+1+1=3.
\]
So
\[
\boxed{15-3=12}
\]
repeated old-token projection rows are avoided.

\subsection*{Attention interactions}
A full-prefix attention call forms a $t\times t$ score matrix, so the three no-cache calls contain
\[
4^2+5^2+6^2=16+25+36=77
\]
query-key interactions.

Cached decoding forms only the newest query against the $t$ available keys, so the same calls contain
\[
4+5+6=15
\]
query-key interactions.

Holding model width fixed, this gives the per-call sequence-length comparison
\[
\boxed{\text{full-prefix no cache: }O(t^2)}
\]
versus
\[
\boxed{\text{cached newest-token attention: }O(t)}.
\]

Cached decoding is not $O(1)$ because the newest query must still compare with an increasingly long K/V history.

\section*{Part F --- Cache size}

For this one-layer, one-head, batch-size-1 toy example,
\[
K_{\mathrm{cache}}\in\mathbb{R}^{L\times d_k},\qquad
V_{\mathrm{cache}}\in\mathbb{R}^{L\times d_v}.
\]
The cache therefore stores
\[
Ld_k+Ld_v=L(d_k+d_v)
\]
scalar elements. Since $d_k=d_v=2$,
\[
\boxed{\text{cache elements}=4L}.
\]
Doubling $L$ doubles the cache size.

\section*{Final conceptual answer}

KV caching is memoization of past K/V projections. It preserves the newest-token attention output because the cached past keys and values are numerically identical to the rows a full-prefix recomputation would regenerate. The cache trades additional memory for less repeated projection work and a smaller attention calculation during each decode step.

\section*{Numerical verification}

The companion \texttt{answer\_key.py} computes the full-prefix causal reference and the cached decode path, then asserts equality of the newest attention row and newest output.

\end{document}
